\section*{参考与延伸阅读}
\addcontentsline{toc}{section}{参考与延伸阅读}

\textbf{讲解路线参考（用户提供的中文字幕资料）}

\begin{enumerate}[leftmargin=2.5em]
\item 《01 从函数到神经网络》
\item 《02 计算神经网络的参数》
\item 《03 调教神经网络的方法》
\item 《04 从矩阵到 CNN》
\item 《05 从词嵌入到 RNN》
\item 《06 简单而强大的 Transformer》
\item 《07 速览大模型 100 词》
\end{enumerate}

这些字幕帮助本书确定了概念引入的\textbf{教学顺序}：从函数与参数开始，经过训练与泛化、CNN、Embedding、RNN，最后进入 Transformer 及大模型应用。书中的数值算例、推导过程、练习与编排为独立重新撰写，不是字幕逐句整理；个别口语化表述已根据数学定义加以辨析。

\textbf{联想记忆与 Transformer 的数学视角}

\begin{itemize}
\item Alex Wang, \emph{You could have invented transformers from memory}:\\
\url{https://alexwang.ai/posts/associative-memory}
\item Vaswani et al., \emph{Attention Is All You Need}, 2017：\\
\url{https://arxiv.org/abs/1706.03762}
\item Goodfellow, Bengio, Courville, \emph{Deep Learning}（公开在线教材）：\\
\url{https://www.deeplearningbook.org/}
\end{itemize}

\textbf{阅读约定}

本文中 $=$ 是数学等式，$\approx$ 通常表示数值近似；讨论高维随机模型时出现的 $\sim$ 往往表示随维度或样本数变化的\textbf{典型数量级/尺度关系}，不必是严格的数值等式。在严谨的渐近分析中，$f(n)\sim g(n)$ 常特指二者比值趋于 $1$；若要表达相差常数因子而增长阶相同，可使用 $\Theta$ 记号。

本书所有国家、词向量、卷积像素和函数拟合样例均为教学用的人工构造数据，不能当作实际模型的已学习参数或真实性能测试。涉及产品、许可与工具生态的内容具有时效性，具体使用应以正式文档为准。
